% ch07_u05 —— 原书 p.276–p.285（7.8.7 -- 7.10.1 开头）
% 本单元自 p.275 末句「$\Delta_{\mathcal{L}}:\mathcal{M}(\mathfrak{g},K)'\to$」的断处续起，
% 至 7.10.1 的「Set」止；尾部接 p.286。

\SourcePage{276}

$\mathcal{M}^{\ell}(K\backslash Y)$ 把 $(\mathfrak{g},K)'$-模 $V$ 送到 $K$-等变左
$\mathcal{D}_{Y}$-模 $(\mathcal{D}_{Y}\otimes\mathcal{L})\otimes_{U(\mathfrak{g}')}V$。
更一般地，
\[
H^{i}_{\mathcal{D}}L\Delta_{\mathcal{L}}(V)=H_{-i}(\mathfrak{g},\mathcal{D}_{Y}\otimes\mathcal{L}\otimes V).
\]

\bdp{7.8.7} \emph{注.} (i) 上述构造只用到 $(\mathfrak{g}',K')$ 在 $(Y,\mathcal{L})$ 上的作用
（不需要整个 $G'$-作用）。
% [Tutor Note: 原书 p.276 此处 "construction" 误拼作「construcion」，底本如此，译文无法保留拼写，照常翻译。]

(ii) 可以证明，$L\Delta_{\mathcal{L}}$ 是 $\Delta_{\mathcal{L}}$ 的左导出函子
（见 7.8.1 的注）。

(iii) 设 $(\mathfrak{g}',K')$ 是 $(\mathfrak{g},K)$ 的平凡扩张，于是
$(\mathfrak{g},K)'$-模与 $(\mathfrak{g},K)$-模相同，而 $\mathcal{L}$ 是
$\mathcal{O}_{Y}$ 配上 $(\mathfrak{g}',K')$ 的显然作用。则
\[
\Delta_{\mathcal{L}}(V)=\mathcal{D}_{Y}\otimes_{U(\mathfrak{g})}V,
\]
即 $\Delta_{\mathcal{L}}$ 与 1.2.4 的函子 $\Delta$ 一致。

\bdp{7.8.8} \emph{命题.} 函子 $L\Delta_{\mathcal{L}}:D(\mathfrak{g},K)'\to D(K\backslash Y)$
是 $\mathcal{H}$-模的态射。

\emph{证明.} 只需证明函子
$\Delta_{\Omega\mathcal{L}}:C(\mathfrak{k}_{\Omega}\times\mathfrak{g},K)'\to C(K\backslash Y,\Omega)$
是 $\mathcal{H}^{c}$-模的态射。

取 7.8.2 中的 $F$、$V$。我们要定义与结合约束相容的 $\Omega$-复形的典范等同
\[
\alpha:\Delta_{\Omega}(F\overset{c}{\circledast}V)\xrightarrow{\ \sim\ }F\overset{c}{\circledast}\Delta_{\Omega}(V).
\]
我们将建立一个与 $K_{\Omega}$-作用相容的典范同构
\[
\tilde{\alpha}:\Delta_{\Omega}(F\widetilde{\circledast}V)\xrightarrow{\ \sim\ }F\widetilde{\circledast}\Delta_{\Omega}(V)
\]
（记号见 7.6.6、7.8.2）。取 $K_{\Omega}$-不变量即得 $\alpha$。

设 $m,p:G\times Y\to Y$ 为作用映射与投影映射，$i:G\times Y\to G\times Y$
为对称态射 $i(g,x)=(g,gx)$；有 $pi=m$。$\mathcal{L}$ 上的 $G'$-作用给出线丛的
$i$-同构 $\tilde{i}:\mathcal{O}_{G}\boxtimes\mathcal{L}_{Y}\xrightarrow{\sim}\mathcal{L}_{G}\boxtimes\mathcal{L}_{Y}$。

下面，对 $\mathfrak{g}$-复形 $P$，记 $C(P)$ 为系数在 $P$ 中、平移 $\dim K$
的李代数链 Chevalley 复形。于是 $C(P)^{\cdot}=C^{\cdot}\otimes P^{\cdot}$，
其中 $C^{a}:=\Lambda^{\dim K-a}\mathfrak{g}$。考虑 $\Omega$-复形
\[
F_{G}\boxtimes\Delta_{\Omega}(V)=F_{G}\boxtimes C(\omega_{Y}\otimes\mathcal{L}_{Y}\otimes V)
\]
与
\[
C((F_{G}\otimes V_{G})\boxtimes(\omega_{Y}\otimes\mathcal{L}_{Y}))
=C((F_{G}\otimes(\mathcal{L}_{G}\otimes V))\boxtimes(\omega_{Y}\otimes\mathcal{L}_{Y}));
\]
这里 $(F_{G}\otimes V_{G})\boxtimes(\omega_{Y}\otimes\mathcal{L}_{Y})$ 上的
$\mathfrak{g}$-作用，是 $\mathfrak{g}'$-作用 $\natural$ 与
$\omega_{Y}\otimes\mathcal{L}_{Y}$ 上的标准 $\mathfrak{g}'$-作用之张量积（见 7.8.2）。

\SourcePage{277}

存在 $\Omega$-复形的典范 $i$-同构
\[
\tilde{\alpha}':F_{G}\boxtimes\Delta_{\Omega}(V)
\xrightarrow{\ \sim\ }C((F_{G}\otimes V_{G})\boxtimes(\omega_{Y}\otimes\mathcal{L}_{Y})),
\]
其定义如下。对 $f\in F_{G}$、$\lambda\in C^{\cdot}$、
$l\in\omega_{Y}\otimes\mathcal{L}_{Y}$、$v\in V$，有
\[
\tilde{\alpha}'(f\otimes\lambda\otimes l\otimes v)
=a(\lambda)\otimes f\otimes\tilde{i}(l)\otimes v;
\]
这里 $a(\lambda)\in\mathcal{O}_{G\times Y}\otimes C^{\cdot}$ 是函数
$a(\lambda)(g,y)=\operatorname{Ad}_{g}(\lambda)$。
$\alpha$ 与微分交换，这一点留给读者验证（用 7.8.2 的注）。

现在有显然的等同
$m\cdot(F_{G}\boxtimes\Delta_{\Omega}(V))=F\widetilde{\circledast}\Delta_{\Omega}(V)$
与
$p\cdot C((F_{G}\otimes V_{G})\boxtimes(\omega_{Y}\otimes\mathcal{L}_{Y}))
=\Delta_{\Omega}(F\widetilde{\circledast}V)$。
于是 $\tilde{\alpha}'$ 给出所要的典范同构 $\tilde{\alpha}$。
$\tilde{\alpha}$ 与 $K_{\Omega}$-作用及结合约束的相容性，留给读者验证。\hfill$\square$

\bdp{7.8.9} 考虑 $Y=G$、$G$ 以左平移作用的情形，并取
$\mathcal{L}=\mathcal{L}_{Y}$ 为 $\mathcal{L}_{G}$ 的对偶线丛（见 7.8.2），
配上左平移给出的显然 $G'$-作用。右 $G'$-平移作用于我们的数据。
因此 $\Omega$-复形 $\Delta_{\Omega}(V)$ 关于 $G'$ 的右平移作用是弱 $G'$-等变的。

设 $C(K\backslash G/^{/}G,\Omega)'\subset C(K\backslash G/^{/}G',\Omega)$
为下述弱 $G'$-等变 $\Omega$-复形 $T$ 构成的子范畴：$\mathbb{G}_{m}\subset G'$
按标准特征标作用于 $T$。记 $D(K\backslash G/^{/}G)'$ 为相应的
$\mathcal{D}$-导出范畴。复形 $\Delta_{\Omega}(V)$ 落在该子范畴中，故有三角函子
$L\Delta:D(\mathfrak{g},K)'\to D(K\backslash G/^{/}G)'$。这些范畴都是
$\mathcal{H}$-模（后者见 7.6.8(v)、7.6.9）。由 7.8.8，$L\Delta$ 是
$\mathcal{H}$-模的态射。7.7.6 与 7.7.11 的一个变体表明，$L\Delta$ 是
t-范畴的等价。
% [Tutor Note: 原书 p.277 此处作「This categories are $\mathcal{H}$-Modules」，主谓数不一致；底本如此，照译为「这些范畴都是」。]

\bdp{7.8.10} \emph{注.} i) 若 $G'$ 是 $G$ 的平凡扩张，则
$D(\mathfrak{g},K)'=D(\mathfrak{g},K)$，且 $L\Delta$ 与 7.7.2 中函子
$\gamma^{-1}$ 所定义的等价一致。这表明 7.7.4 与 7.8.3 中的 Hecke 作用确实一致。

(ii) 设扩张是任意的。则拉回函子
$r:D(K\backslash G/K)\to D(K'\backslash G'/K')$ 是幺半范畴的态射，

\SourcePage{278}

而全忠实嵌入 $D(\mathfrak{g},K)'\hookrightarrow D(\mathfrak{g}',K')$ 是
$\mathcal{H}$-模的 $r$-态射。于是在 $(\mathfrak{g}',K')$ 的情形，扭的图景本质上
等价于非扭的图景。不过在应用中，方便的做法是把扭（又称水平，又称中心荷）
单独保留。
% [Tutor Note: 原书 p.278 此处作「is $r$-Morphism of Hecke Modules」，英文语法不规范（应为 a morphism）；底本如此，照译为「是 $\mathcal{H}$-模的 $r$-态射」。]

\bdp{7.8.11} 下面解释 7.1.1 中「Hecke 模式」(d) 的 $\Gamma$ 部分。
本文主体不需要这一题材，读者可以跳过本节余下部分。我们处理扭的版本，
故处于 7.8.6 的情形。对 $T\in C(K\backslash Y,\Omega)$，$Y$ 上的
$\mathcal{D}$-复形 $\mathcal{D}T_{Y}$ 是 $K$-等变的（见 7.6.11）。
把 $\mathcal{D}T_{Y}$ 看作带 $(\mathfrak{k}_{\Omega}\times\mathfrak{g},K)$-作用的
$\mathcal{O}$-复形。置
\[
\Gamma_{\mathcal{L}}(T):=\Gamma(Y,\mathcal{D}T_{Y}\otimes(\omega_{Y}\otimes\mathcal{L}_{Y})^{*}).
\]
这是 Harish-Chandra 复形（回忆 $(\mathfrak{g}',K)$ 作用于
$\omega_{Y}\otimes\mathcal{L}_{Y}$），于是有 DG 函子
$\Gamma_{\mathcal{L}}:C(K\backslash Y,\Omega)\to C(\mathfrak{k}_{\Omega}\times\mathfrak{g},K)'$。
设
\[
R\Gamma_{\mathcal{L}}:D(K\backslash Y)\to D(\mathfrak{g},K)'
\]
为其右导出函子。若 $T$ 是松的，则
$\Gamma_{\mathcal{L}}(T)=R\Gamma_{\mathcal{L}}(T)$，故 $R\Gamma_{\mathcal{L}}$
的定义是正当的。

注意 $R\Gamma_{\mathcal{L}}$ 是左 t-正合函子；记
$\Gamma_{\mathcal{L}}:\mathcal{M}(K\backslash Y)\to\mathcal{M}(\mathfrak{g},K)'$
为相应的左正合函子。有
$\Gamma_{\mathcal{L}}(M)=\Gamma(Y,M\otimes(\omega_{Y}\otimes\mathcal{L}_{Y})^{*})$。
若处于 7.8.7(iii) 的情形，则在右 $\mathcal{D}$-模与左 $\mathcal{D}$-模的标准
等同之下，该函子与 1.2.4 的函子 $\Gamma$ 一致。

\bdp{7.8.12} \emph{引理.} 函子 $R\Gamma_{\mathcal{L}}$ 是 $\mathcal{H}$-模的态射。

\emph{证明.} 只需证明 $\Gamma_{\mathcal{L}}$ 是 $\mathcal{H}^{c}$-模的态射，
即对 $F\in\mathcal{H}^{c}$ 与上述 $T$，定义与结合约束相容的典范同构
$\beta:\Gamma_{\mathcal{L}}(F\overset{c}{\circledast}T)\xrightarrow{\sim}
F\overset{c}{\circledast}\Gamma_{\mathcal{L}}(T)$。我们先写出与
$K_{\Omega}$-作用相容的典范同构
$\tilde{\beta}:\Gamma_{\mathcal{L}}(F\widetilde{\circledast}T)
\xrightarrow{\sim}F\widetilde{\circledast}\Gamma_{\mathcal{L}}(T)$；
取 $K_{\Omega}$-不变量即得 $\beta$。

$\mathcal{L}$ 上的 $G'$-作用给出同构
$m_{Y}^{*}((\omega_{Y}\otimes\mathcal{L}_{Y})^{*})
=\mathcal{L}_{G}\otimes(\omega_{Y}\otimes\mathcal{L}_{Y})^{*}$，
而 $\mathcal{D}_{Y}$ 上的 $G$-作用（作为左 $\mathcal{O}_{Y}$-模）给出
% [Tutor Note: 原书 p.278 末行「(as on a left $\mathcal{O}_Y$-module」的右括号缺失；底本如此，译文补足了右括号。]

\SourcePage{279}

同构 $m_{Y}^{*}(\mathcal{D}_{Y})=\mathcal{O}_{G}\boxtimes\mathcal{D}_{Y}$。
这些同构把 $\Gamma_{\mathcal{L}}(F\widetilde{\circledast}T)^{\cdot}$ 与
\[
\Gamma(G\times Y,(F^{\cdot}\otimes\mathcal{L}_{G})\boxtimes
(\mathcal{D}T_{Y}^{\cdot}\otimes(\omega_{Y}\otimes\mathcal{L}_{Y})^{*}))
\]
等同起来。该向量空间等于
$\Gamma(G,F^{\cdot}\otimes\mathcal{L}_{G})\otimes
\Gamma(Y,\mathcal{D}T_{Y}^{\cdot}\otimes(\omega_{Y}\otimes\mathcal{L}_{Y})^{*})$，
即 $(F\widetilde{\circledast}\Gamma_{\mathcal{L}}(T))^{\cdot}$。
我们的 $\tilde{\beta}$ 就是这些等同的复合。$\tilde{\beta}$ 是 Harish-Chandra
复形的同构且与 $K_{\Omega}$-作用相容，这一点留给读者验证。\hfill$\square$

\bdsec{7.9. 额外的对称性与参数}

\bdp{7.9.1} 在本文主体（即 5.4）中，我们用到 7.1.1 的 Hecke 模式的一个
等变版本。确切地说，给定一个额外的 Harish-Chandra 对 $(\mathfrak{l},P)$，
它作用于 $(G,K)$，而我们要找 7.1.1(a)--(d) 的 $(\mathfrak{l},P)$-等变版本。
下面简要说明这一设定；全部细节见本节余下部分。Hecke 范畴 $\mathcal{H}$ 是
$K\backslash G/K$ 上弱 $(\mathfrak{l},P)$-等变 $\mathcal{D}$-模范畴的导出版本。
这是幺半三角范畴（即 7.1.1(a) 在当前设定下的类似物）。$\mathcal{H}$ 作用于
$(\mathfrak{l}\bowtie\mathfrak{g},P\bowtie K)$-模的适当导出范畴 $D_{HC}$；
这是类似于 7.1.1(c) 的 Harish-Chandra 对应物。几何对应物如下。设 $X$ 是带
$(\mathfrak{l},P)$-结构 $X^{\wedge}$ 的「参数」空间（见 2.6.4）。考虑以
$X^{\wedge}$ 为参数、带 $G$-作用的光滑簇族 $Y^{\wedge}$。设 $X^{\wedge}$ 上的
$(\mathfrak{l},P)$-作用提升到 $Y^{\wedge}$，且与 $G$-作用相容。则 $\mathcal{H}$
作用于 $X$-栈 $\mathcal{B}=(P\bowtie K)\backslash Y^{\wedge}$ 的
$\mathcal{D}$-模导出范畴 $D(\mathcal{B})$（这是 7.1.1(b) 的版本）。
有适当的局部化函子 $L\Delta:D_{HC}\to D(\mathcal{B})$，它与 Hecke 作用交换
（这就是 7.1.1(d)）。对带 $(\mathfrak{l},P)$-作用的代数 $A$，上述构造有
$A$-线性版本：考察带 $A$-作用的 Harish-Chandra 模与带 $A_{X}$-作用的
$\mathcal{D}$-模（$A_{X}$ 的定义见 2.6.6）。相应的三角范畴记作
$\mathcal{H}_{A}$、$D_{HCA}$ 与 $D(\mathcal{B},A_{X})$。

\SourcePage{280}

这些构造本质上是前几节构造的直截了当的修改；我们写下它们，是为在 5.4 中
直接引用。

\emph{注.} 等变 Hecke 模式不能归约为把 $G$ 换成对应于 Harish-Chandra 对
$(\mathfrak{l}\bowtie\mathfrak{g},P\bowtie G)$ 的群归纳概形后的普通模式。
事实上，我们的 $\mathcal{H}$ 比相应的「普通」Hecke 范畴大得多：后者由
$K\backslash G/K$ 上强 $P$-等变 $\mathcal{D}$-模构成。特别地，$\mathcal{H}$
以 $(\mathfrak{l},P)$-模的张量范畴为张量子范畴。上述纤维化结构 $Y/X$ 正是
以使整个 $\mathcal{H}$ 作用于 $D(\mathcal{B})$ 所需。
% [Tutor Note: 原书 p.280 此处作「much larger then」，应为「than」；底本如此，照译为「大得多」。]

\bdp{7.9.2} 于是考虑作用于 $(G,K)$ 的 Harish-Chandra 对 $(\mathfrak{l},P)$。
这里 $P$ 可以是任意仿射群概形（不必有限型），但我们假设
$\operatorname{Lie}P$ 在 $\mathfrak{l}$ 中余维有限。考虑 DG 范畴
$\mathcal{H}^{c}$，其对象是 $K\backslash G/K$ 上的 $\Omega$-复形 $F$，
带 $(\mathfrak{l},P)$-作用，且该作用提升 $G/K$ 上的 $(\mathfrak{l},P)$-作用。
这样的 $F$ 等同于 $G$ 上
$(\mathfrak{l},P)\bowtie(K_{\Omega}\times K_{\Omega})$-等变的 $\Omega$-复形。
称 $\mathcal{H}^{c}$ 为 $(\mathfrak{l},P)$-等变预 Hecke 范畴。
$\mathcal{H}^{c}$ 的同伦范畴中那些作为普通 $\Omega$-复形的
$\mathcal{D}$-拟同构的态射构成局部化族。$(\mathfrak{l},P)$-等变 Hecke 范畴
$\mathcal{H}$ 就是相应的局部化。于是 $\mathcal{H}$ 是 t-范畴，其核心等于
$G/K$ 上带弱 $(\mathfrak{l}\bowtie\mathfrak{k},P\bowtie K)$-作用的
$\mathcal{D}$-模范畴（这里 $K$ 以左平移作用于 $G/K$），且 $K$ 的作用实际上是
强的。

现在 $\mathcal{H}^{c}$ 是 DG 幺半范畴，$\mathcal{H}$ 是幺半三角范畴。
事实上，7.6.1 的全部定义在当前情形都适用。

\emph{注.} 取 Harish-Chandra 模 $V\in\mathcal{M}(\mathfrak{l},P)$。
把它对应到 $G/K$ 的指定点处的摩天楼层，看作集中在 $0$ 次、带平凡
$K_{\Omega}$-作用的 $\Omega$-复形。这是 $\mathcal{H}^{c}$ 的对象。函子
$\mathcal{M}(\mathfrak{l},P)\to\mathcal{H}^{c},\mathcal{H}$ 是全忠实幺半函子。
注意 $\mathcal{M}(\mathfrak{l},P)$ 以典范方式属于（预）Hecke 幺半范畴的中心，
即对任意上述 $V$、

\SourcePage{281}

$F\in\mathcal{H}$，有与 $F$ 的张量积及 $V$ 的张量积都相容的典范同构
$V\circledast{}F\xrightarrow{\sim}F\circledast{}V$。事实上，两个对象都等于
$V\otimes F$。

\bdp{7.9.3} 为定义 $\mathcal{D}$-模上的 Hecke 作用，需要先作些准备。

设 $X$ 是光滑簇，$Y$ 是 $\mathcal{D}_{X}$-概形。$Y$ 上的
$\mathcal{D}_{X}\Omega_{/X}$-复形是带 $\mathcal{D}_{X}$-结构的
DG $\Omega_{Y/X}$-模（$\mathcal{D}_{X}$-结构 $:=$ 沿 $Y/X$ 上结构联络的叶的
平坦联络）。确切地说，$Y$ 上的 $\mathcal{D}_{X}$-结构在
$\Omega_{Y/X}(\mathcal{D}_{X}):=\mathcal{D}_{X}\otimes_{\mathcal{O}_{X}}\Omega_{Y/X}$
上定义了结合 DG 代数结构。于是 $Y$ 上的 $\mathcal{D}_{X}\Omega_{/X}$-复形是
左 DG $\Omega_{Y/X}(\mathcal{D}_{X})$-模，且作为 $\mathcal{O}_{Y}$-模是拟凝聚的。

$Y$ 上 $\mathcal{D}_{X}\Omega_{/X}$-复形的 DG 范畴
$C(Y,\mathcal{D}_{X}\Omega_{/X})$ 是张量范畴（张量积取在 $\Omega_{Y/X}$ 上）。
拉回函子
\[
C(\mathcal{M}^{\ell}(X))\longrightarrow C(Y,\mathcal{D}_{X}\Omega_{/X}),
\qquad
M\mapsto\Omega_{Y/X}\otimes_{\mathcal{O}_{X}}\Omega_{Y/X},
\]
是张量函子。
% [Tutor Note: 原书 p.281 此处的拉回函子印作 $M\mapsto\Omega_{Y/X}\otimes_{\mathcal{O}_X}\Omega_{Y/X}$，第二个因子应为 $M$；底本如此，保留原样。]
特别地 $C(Y,\mathcal{D}_{X}\Omega_{/X})$ 是 $\mathcal{M}^{\ell}(X)$-模
（有 $M\circledast F=M\otimes_{\mathcal{O}_{X}}F$）。

注意，对 $Y$ 上的 $\mathcal{D}_{X}\Omega_{/X}$-复形 $F$，有绝对 $\Omega$-复形
$\Omega_{X}F$，它定义为沿 $X$ 的 de Rham 复形，系数在 $F$ 中
\footnote{如同 7.2，函子 $\Omega_{X}$ 有左伴随函子 $\mathcal{D}_{X}$。}。
于是若 $Y$ 是光滑簇，就有 $\mathcal{D}_{X}\Omega_{/X}$-复形的
$\mathcal{D}$-拟同构概念。$C(Y,\mathcal{D}_{X}\Omega_{/X})$ 的同伦范畴相应的
局部化记作 $D(Y,\mathcal{D}_{X}\Omega_{/X})$。函子
$\Omega_{X}:D(Y,\mathcal{D}_{X}\Omega_{/X})\longrightarrow D(Y,\Omega)$
是三角范畴的等价。

\bdp{7.9.4} 现在设 $X$ 是带 $(\mathfrak{l},P)$-结构 $X^{\wedge}$ 的光滑簇
（见 2.6.4）。设 $Y^{\wedge}$ 是带 $(\mathfrak{l},P)\bowtie G$-作用的概形，
且有与作用相容的光滑态射 $p^{\wedge}:Y^{\wedge}\to X^{\wedge}$（于是 $G$ 沿
纤维作用，$p^{\wedge}$ 与 $(\mathfrak{l},P)$ 的作用交换）。置
$Y:=P\backslash Y^{\wedge}$。这是带光滑投影 $p:Y\to X$ 的光滑簇。
$Y^{\wedge}$ 上的 $(\mathfrak{l},P)$-作用在 $Y$ 上定义了 $\mathcal{D}_{X}$-概形
结构。$Y^{\wedge}$ 上的 $G$-作用在 $Y$ 上给出水平的 $G_{X}$-作用（群
$\mathcal{D}_{X}$-概形 $G_{X}$ 的定义见 2.6.6）。
% [Tutor Note: 原书 p.281 此处 "horizontal" 误拼作「horisontal」（p.283 同）；底本如此，照译为「水平的」。]

\SourcePage{282}

考虑 $X$ 上纤维化的栈 $\mathcal{B}:=K_{X}\backslash Y=(P\bowtie K)\backslash Y^{\wedge}$，
于是有相应的左 $\mathcal{D}$-模范畴 $\mathcal{M}^{\ell}(\mathcal{B})$ 与
$\mathcal{B}$ 上 $\Omega$-复形的 t-范畴 $D(\mathcal{B})$。该 t-范畴还有另一种
用 $\mathcal{D}_{X}\Omega_{/X}$-复形的实现，下面描述。

考虑 DG 群 $\mathcal{D}_{X}$-概形
$G_{\Omega X}:=(G_{X},\Omega_{G_{X}/X})$、$K_{\Omega X}$。
$Y$ 上 $\mathcal{D}_{X}\Omega_{/X}$-复形的 $K_{\Omega X}$-作用按 7.6.4 定义。
现在有 $Y$ 上 $K_{\Omega X}$-等变 $\mathcal{D}_{X}\Omega_{/X}$-复形的 DG 范畴
$C(K_{X}\backslash Y,\mathcal{D}_{X}\Omega_{/X})$。把它的同伦范畴按
$\mathcal{D}$-拟同构局部化，得到三角范畴
$D(K_{X}\backslash Y,\mathcal{D}_{X}\Omega_{/X})$。de Rham 函子 $\Omega_{X}$
把它与 $D(\mathcal{B})$ 等同。

现在可以定义 $D(\mathcal{B})$ 上的 Hecke 作用。先构造 $\mathcal{H}^{c}$ 在
$C(K_{X}\backslash Y,\mathcal{D}_{X}\Omega_{/X})$ 上的作用
$\overset{c}{\circledast}$。事实上，对 $F\in\mathcal{H}^{c}$，有 $G_{X}$ 上的
$\mathcal{D}_{X}\Omega_{/X}$-复形 $F_{X}$，它关于左、右平移是
$K_{\Omega X}$-等变的。于是对
$T\in C(K_{X}\backslash Y,\mathcal{D}_{X}\Omega_{/X})$，在 $\mathcal{D}_{X}$-概形
$G_{X}\times Y$（$G_{X}$ 与 $Y$ 在 $X$ 上的纤维积）上有
$\mathcal{D}_{X}\Omega_{/X}$-复形 $F\boxtimes T$。它关于 $G_{X}\times Y$ 上所有
$K_{X}$-作用是 $K_{\Omega X}$-等变的。故 $F\boxtimes T$ 下降到
$G_{X}\times_{K_{X}}Y$。定义
$F\overset{c}{\circledast}T\in C(K_{X}\backslash Y,\mathcal{D}_{X}\Omega_{/X})$
为上述复形沿作用映射 $G_{X}\times_{K_{X}}Y\to Y$ 的前推。Hecke 作用
$\circledast:H\times D(\mathcal{B})\longrightarrow D(\mathcal{B})$ 是
$\overset{c}{\circledast}$ 的右导出函子；按惯例，可以用松的
$\mathcal{D}_{X}\Omega_{/X}$-复形来计算它。

\emph{注.} 对 $W\in\mathcal{M}(\mathfrak{l},P)\subset\mathcal{H}^{c}$ 与上述
$T$，有 $W\overset{c}{\circledast}T=W\circledast{}T=W_{X}\otimes T$
（$\mathcal{D}_{X}$-模 $W_{X}$ 的定义见 2.6.6）。

\bdp{7.9.5} 下面定义 Harish-Chandra 范畴。设 $G'$ 如 7.8.1，并设给定
$G$ 上 $(\mathfrak{l},P)$-作用到 $G'$ 上的提升，它保持 $K\subset G'$ 且固定
$\mathbb{G}_{m}\subset G'$。于是有 Harish-Chandra 对
$(\mathfrak{l},P)\bowtie(\mathfrak{g}',K')$。设 $C_{HC}$ 为
$(\mathfrak{l},P)\bowtie(\mathfrak{k}_{\Omega}\times\mathfrak{g},K)'$-复形的
范畴，即带相容 $(\mathfrak{l},P)$-作用的
$(\mathfrak{k}_{\Omega}\times\mathfrak{g},K)'$-复形（记号见 7.8.1）。设
$D_{HC}$ 为相应的导出范畴。这是 t-范畴，核心为
$\mathcal{M}_{HC}=\mathcal{M}(\mathfrak{l}\bowtie\mathfrak{g},P\bowtie K)'$。
下面

\SourcePage{283}

我们把 $C_{HC}$、$D_{HC}$ 的对象简称为 Harish-Chandra 复形，把
$\mathcal{M}_{HC}$ 的对象简称为 Harish-Chandra 模。

预 Hecke 范畴 $\mathcal{H}^{c}$ 作用于 $C_{HC}$。事实上，7.8.2 的构造在
当前情形完全说得通（$(\mathfrak{l},P)$ 通过结构搬运作用于
$F\overset{c}{\circledast}V$）。$D_{HC}$ 上的 $\mathcal{H}$-作用
$\circledast$ 是 $\overset{c}{\circledast}$ 的右导出函子。7.8.4--7.8.5 的
结果原样移植到当前设定。

\emph{注.} 对 $W\in\mathcal{M}(\mathfrak{l},P)\subset\mathcal{H}^{c}$ 与
Harish-Chandra 复形 $V$，有 Harish-Chandra 复形的典范同构
$W\overset{c}{\circledast}V=W\circledast{}V=W\otimes V$。

\bdp{7.9.6} 下面转向局部化函子。7.8.6 的构造移植到我们的设定如下。
从 7.9.4 的 $Y^{\wedge}$ 出发。设它带线丛 $\mathcal{L}_{Y^{\wedge}}$，且
$Y^{\wedge}$ 上的 $(\mathfrak{l},P)\bowtie G$-作用提升为
$\mathcal{L}_{Y^{\wedge}}$ 上的 $(\mathfrak{l},P)\bowtie G'$-作用，使
$\mathbb{G}_{m}\subset G'$ 按与标准特征标相反的特征标作用。设
$\mathcal{L}_{Y}$ 是 $\mathcal{L}_{Y^{\wedge}}$ 沿 $P$-作用下降到 $Y$ 的线丛。
该线丛带典范 $\mathcal{D}_{X}$-结构，它来自 $\mathcal{L}_{Y^{\wedge}}$ 上的
$\mathfrak{l}$-作用。它还带水平的 $G'_{X}$-作用。

有 DG 函子
\begin{equation}\tag{338}
\Delta_{\Omega}=\Delta_{\Omega\mathcal{L}}:C_{HC}\longrightarrow
C(K_{X}\backslash Y,\mathcal{D}_{X}\Omega_{/X}),
\end{equation}
$\Delta_{\Omega}(V)=\operatorname{Hom}_{\mathfrak{g}_{X}}
(DR_{\mathfrak{g}_{X}},\omega_{Y/X}\times\mathcal{L}_{Y}\times V)[-\dim K]$
（对照（336））。
% [Tutor Note: 原书 p.283 此处印作 $\omega_{Y/X}\times\mathcal{L}_Y\times V$（用 $\times$），而 7.8.6 的（336）式用 $\otimes$；底本如此，保留原样。]
如 7.8.6，该函子把拟同构送到 $\mathcal{D}$-拟同构，故给出三角函子
\begin{equation}\tag{339}
L\Delta=L\Delta_{\mathcal{L}}:D_{HC}\longrightarrow D(\mathcal{B}),
\end{equation}
它是右 t-正合的。两个核心之间的相应右正合函子
$\Delta_{\mathcal{L}}:\mathcal{M}_{HC}\longrightarrow\mathcal{M}^{\ell}(\mathcal{B})$
把 $V$ 送到 $K_{X}$-等变左 $\mathcal{D}_{Y}$-模
\[
(\mathcal{D}_{Y/X}\otimes\mathcal{L}_{Y})\otimes_{U(\mathfrak{g}'_{X})}V_{X}.
\]

函子 $\Delta_{\Omega}$、$L\Delta$ 与 Hecke 作用交换。事实上，7.8.8 的证明
逐字照搬到我们的设定。特别地，对任意 $W\in\mathcal{M}(\mathfrak{l},P)$、
$V\in D_{HC}$，有 $L\Delta(W\otimes V)=W_{X}\otimes L\Delta(V)$。

\SourcePage{284}

\bdp{7.9.7} \emph{$A$-线性版本.} 设此外还给定带 $(\mathfrak{l},P)$-作用的
交换代数 $A$。把它接到上述模式上如下。

(i) 记 $\mathcal{H}^{c}_{A}$ 为满足如下条件的对象 $F\in\mathcal{H}^{c}$ 的
DG 范畴：$F$ 带 $A$-作用，$A$ 与 $(\mathfrak{l},P)$ 的作用相容，且 $F$ 是
$A$-平坦的。设 $\mathcal{H}_{A}$ 为相应的 $\mathcal{D}$-导出范畴。卷积积按
7.9.2 定义（张量积取在 $A$ 上），故 $\mathcal{H}^{c}_{A}$ 与
$\mathcal{H}_{A}$ 是幺半范畴。设 $\mathcal{M}(\mathfrak{l},P)^{fl}_{A}$ 为
带 $(\mathfrak{l},P)$-作用的平坦 $A$-模的张量范畴。如 7.9.2 的注，有典范
全忠实幺半函子
$\mathcal{M}(\mathfrak{l},P)^{fl}_{A}\longrightarrow\mathcal{H}^{c}_{A},\mathcal{H}_{A}$，
它们把 $\mathcal{M}(\mathfrak{l},P)^{fl}_{A}$ 送到 Hecke 范畴的中心。

(ii) 设处于 7.9.4 的情形。考虑 $\mathcal{B}$ 上带 $A_{X}$-作用的左
$\mathcal{D}$-模范畴 $\mathcal{M}^{\ell}(\mathcal{B},A_{X})$（$\mathcal{D}_{X}$-代数
$A_{X}$ 的定义见 2.6.6）。设 $C(\mathcal{B},A_{X}\otimes\Omega)$ 为
$\mathcal{B}$ 上带 $A_{X}$-作用的 $\Omega$-复形的 DG 范畴，
$D(\mathcal{B},A_{X})$ 为相应同伦范畴关于 $\mathcal{D}$-拟同构的局部化。
这是 t-范畴，核心为 $\mathcal{M}^{\ell}(\mathcal{B},A_{X})$。如 7.9.4，也可以
用 $\mathcal{D}_{X}\Omega_{/X}$-复形来定义该 t-范畴。确切地说，设
$C(K_{X}\backslash Y,A_{X}\mathcal{D}_{X}\Omega_{/X})$ 为
$C(K_{X}\backslash Y,\mathcal{D}_{X}\Omega_{/X})$ 中带 $A_{X}$-作用（与
$K_{\Omega X}$-作用交换）的对象构成的 DG 范畴。按 $\mathcal{D}$-拟同构
局部化，得到三角范畴
$D(K_{X}\backslash Y,A_{X}\mathcal{D}_{X}\Omega_{/X})$。de Rham 函子
$\Omega_{X}$ 把它与 $D(\mathcal{B},A_{X})$ 等同。

$A$-线性设定下的 Hecke 作用完全按 7.9.4 定义。7.9.4 的注仍然成立
（张量积取在 $A_{X}$ 上）。

(iii) 设处于 7.9.5 的情形。定义 $C_{HCA}$ 为带相容 $A$-作用的
Harish-Chandra 复形范畴（于是 $A$ 与
$(\mathfrak{k}_{\Omega}\times\mathfrak{g},K)'$ 的作用交换）。设 $D_{HCA}$ 为
相应的导出范畴。这是 t-范畴，核心 $\mathcal{M}_{HCA}$ 等于带相容 $A$-作用的
$(\mathfrak{l}\bowtie\mathfrak{g},P\bowtie K)'$-模范畴。所有

\SourcePage{285}

关于 Hecke 作用的构造与结果都原样成立。在 7.9.5 的注中取
$W\in\mathcal{M}(\mathfrak{l},P)^{fl}_{A}$；张量积 $W\otimes V$ 取在 $A$ 上。
局部化函子的 $A$-线性设定无需改动。

\emph{注.} 有显然的函子（用 $A$ 作张量），它们把上述普通范畴送到附上 $A$
的范畴。这些函子与我们考虑的所有结构相容。忘掉 $A$-作用的函子
$D(\mathcal{B},A_{X})\to D(\mathcal{B})$、$D_{HCA}\to D_{HC}$ 是
$\mathcal{H}$-模的态射。它们与局部化函子交换。

\bdp{7.9.8} \emph{变体.} 设除 $A$ 之外还给定与 $(\mathfrak{l},P)$-作用相容的
交换代数态射 $e:\mathfrak{z}\to A$。这里
$\mathfrak{z}:=U(\mathfrak{g})^{\prime\operatorname{Ad}G}$（故若 $G$ 连通，
则 $\mathfrak{z}$ 是 $U(\mathfrak{g})'$ 的中心）。于是 $\mathfrak{z}$ 以两种
方式作用于 $\mathcal{M}_{HCA}$ 或 $C_{HCA}$ 的任一对象。记
$\mathcal{M}^{e}_{HCA}$、$C^{e}_{HCA}$ 为 $\mathfrak{z}$ 的两种作用一致的
对象构成的范畴；设 $D^{e}_{HC}$ 为相应的导出范畴。$\mathcal{H}^{c}_{A}$ 在
$C_{HCA}$ 上的作用是 $\mathfrak{z}$-线性的（见 7.8.4），故它保持
$C^{e}_{HCA}$。于是有 $\mathcal{H}_{A}$ 在 $D^{e}_{HCA}$ 上的作用。显然函子
$D^{e}_{HCA}\to D_{HCA}$ 是 $\mathcal{H}_{A}$-模的态射。

\emph{注.} 若 $e$ 是满射，则 $\mathcal{M}^{e}_{HCA}$ 是 $\mathcal{M}_{HC}$ 中
由被 $\operatorname{Ker}e$ 零化的 Harish-Chandra 模构成的满子范畴。
$C^{e}_{HCA}$ 亦然。

\bdsec{7.10. $\mathcal{D}$-晶体}

下文概述 $\mathcal{D}$-模理论的晶体方法。与通常的形式体系不同，它不区分
光滑与不光滑的概形。

本节中，「概形」指「局部有限型的 $\mathbb{C}$-概形」。代数空间与栈亦然。
形式概形或代数空间假设为局部归纳有限型的
\footnote{即：任一闭子概形都是有限型的。}。

\bdp{7.10.1} 设 $f:Y\to X$ 是概形的拟有限态射。则 Grothendieck 的函子
$Rf^{!}:D^{b}(X,\mathcal{O})\to D^{b}(Y,\mathcal{O})$ 左 t-正合。置

% 尾部接 p.286：$f^{!}:=H^{0}Rf^{!}:\mathcal{M}(X,\mathcal{O})\to\mathcal{M}(Y,\mathcal{O})$，……
